\chapter{実験1}

\section{本章の概要}

本章では、処理分岐分類性能を評価する実験1について述べる。評価目的、比較対象、評価方針、および実験完了後に確認したい観点を整理する。

\section{実験1の設定}

実験1の目的は、提案手法が \texttt{act} / \texttt{reobserve} / \texttt{ask\_user} / \texttt{stop} / \texttt{defer} を適切に分類できるかを評価することである。比較対象は、Zero-shot VLM、VLM + 分岐定義のみ、VLM + 工程条件リストである。評価では、Resolution Mode Accuracy および Macro F1 を主要指標とし、提案手法がベースラインより安定した分類傾向を示すかを確認する。

\section{評価の観点}

本研究は執筆時点で評価進行中であるため、本節では実測結果ではなく、実験完了後に確認すべき観点を述べる。具体的には、提案手法が Zero-shot VLM や分岐定義のみを与えた条件と比べて、処理分岐分類の安定性をどの程度改善するかに着目する。特に、\texttt{ask\_user} と \texttt{reobserve} の混同、\texttt{act} と \texttt{reobserve} の混同がどの程度抑えられるかを確認する予定である。

\section{ベースラインとの比較}

構造化診断を行わない Zero-shot VLM では、画像に現れた曖昧さに引きずられて、現在工程に依存した切り替えが弱くなる可能性がある。また、分岐定義のみを与えた条件では、各分岐の意味は理解できても、どの工程条件が blocking であるかの判断が不安定になると考えられる。これに対し、提案手法は工程条件、証拠充足性、不確実性原因を明示的に扱うため、実験完了後には判断の一貫性がどの程度向上するかを比較できる。

\section{考察}

提案手法が分類精度を改善する傾向を示した場合、VLM に単に分岐ラベルを選ばせるよりも、不確実性を中間概念に分解して診断させることが有効であると解釈できる。一方で、\texttt{stop} や \texttt{ask\_user} を過剰に選択する傾向が確認された場合は、安全性重視と作業継続性のバランスを再検討する必要がある。
